\documentclass[10pt,a4paper]{amsart}
\usepackage[margin=19mm]{geometry}
\usepackage{amsmath,amssymb,mathtools}
\usepackage[scr=rsfs,cal=euler]{mathalfa}
\usepackage{xfrac}
\usepackage[unicode,hidelinks,bookmarks=false]{hyperref}
\newcommand{\ie}{i.e.\ }
\newcommand{\cf}{cf.\ }
\newcommand{\resp}{respectively\ }
\newcommand{\Subsection}{\subsection}
\providecommand{\nobreakdash}{\nobreak\hbox{-}}
\setlength{\emergencystretch}{2em}
\multlinegap0pt
\usepackage{bm}
\usepackage{bbm}
\usepackage{dsfont}
\usepackage{xspace}
\usepackage{cancel}
\usepackage{mathtools}
\usepackage{cleveref}
\usepackage{amsthm}
\makeatletter
\@ifundefined{definition}{\newtheorem{definition}{Definition}}{}
\@ifundefined{lemma}{\newtheorem{lemma}[definition]{Lemma}}{}
\@ifundefined{proposition}{\newtheorem{proposition}[definition]{Proposition}}{}
\@ifundefined{theorem}{\newtheorem{theorem}[definition]{Theorem}}{}
\makeatother
\DeclareMathOperator*{\diam}{diam}
\DeclareMathOperator*{\supp}{supp}
\newcommand{\dist}{\mathrm{dist}}
\renewcommand{\epsilon}{\varepsilon}
\title[Structure of average distance minimizers in general dimensio]{Structure of average distance minimizers in general dimensions}
\author{Lucas O'Brien, Forest Kobayashi,, Young-Heon Kim}
\date{Source: arXiv:2503.23256v3 (CC BY 4.0)}
\begin{document}
\maketitle

\noindent\emph{Source and licence.} Structure of average distance minimizers in general dimensions. Version arXiv:2503.23256v3 (CC BY 4.0), distributed under the Creative Commons Attribution 4.0 International licence, as recorded in the arXiv licence metadata for this exact version and shown on the abstract page https://arxiv.org/abs/2503.23256v3. Full rights notice: licenses/arxiv-2503.23256-CC-BY-4.0.txt.\par\smallskip

\noindent\textit{Editorial scope.} Counted unit: the Theorem whose source label is \texttt{prop:boundingmassofnoncutpoints} in the article (source0.tex lines 1532--1545, source label \texttt{prop:boundingmassofnoncutpoints}), delivered together with the article's own complete original proof of it (source0.tex lines 1547--1732, one \texttt{proof} environment of 186 lines). No mathematics is written, repaired or re-derived: statement and proof are the article's own text line for line. Required context retained uncounted: eq:adp (lines 152--162); lemma:basicinequality (lines 676--683); strongbarycentreapproximation (lines 925--942); def:trivialbarycentrefield (lines 1088--1098); approximatebysmooth (lines 1153--1164); ridgesetnegligable (lines 1338--1343); noncutneighbourhood (lines 1513--1527). Every definition, display and label of those lines is retained unchanged. Omitted from this package and not redistributed: 1--151, 163--675, 684--924, 943--1087, 1099--1152, 1165--1337, 1344--1512, 1528--1531, 1546--1546, 1733--3185. Nothing inside a retained excerpt is omitted or altered: every retained body is delivered exactly as the article's lines are written, so no omission receipt of any kind accompanies this package. Citations: the retained excerpts carry no citation command at all, so no bibliography entry of the article is transcribed and no reference is redistributed. Reference adaptation: none was needed and none was made; every reference inside the delivered unit resolves inside it, so no unresolved reference or citation is produced. Printed slips kept literal: no printed word, symbol, hypothesis or number of the counted statement or of its proof was altered. Layout and class: the article's own class is \texttt{amsart} with packages including amsfonts, amsmath, amsthm, amssymb, amsrefs, mathtools, mathrsfs, mathalfa, enumitem, hyperref, cleveref; the wrapper is the lane's pinned \texttt{amsart} editorial wrapper at 10pt on A4, which restates the theorem-like environments the retained text needs and adds the article's own macro definitions that the retained text needs (\texttt{\textbackslash diam}, \texttt{\textbackslash dist}, \texttt{\textbackslash epsilon}, \texttt{\textbackslash supp}), with extra wrapper packages (bm, bbm, dsfont, xspace, cancel, mathtools, cleveref). The delivered numbering is the unit's own. Assets: no retained span contains a figure environment, a graphics-inclusion command, a PDF-page inclusion, a table or a TikZ picture; no image file is copied into this package, the package declares no asset dependency and no image file is redistributed with it. Licence: version arXiv:2503.23256v3 of this article is distributed under the Creative Commons Attribution 4.0 International licence, as recorded in the arXiv licence metadata for this exact version and shown on the abstract page https://arxiv.org/abs/2503.23256v3. A full rights notice is delivered at licenses/arxiv-2503.23256-CC-BY-4.0.txt. Characters: every retained excerpt is pure ASCII, so the delivered engine, XeLaTeX with its default Latin Modern text font, reports no missing character.
\begin{equation}
  \label{eq:adp}
  \begin{cases}
    \text{Minimize:} \quad
    & \mathscr J_p(\Sigma) \\
    \text{Subject to:} \quad
    & \Sigma \text{ is compact and connected, and } \mathcal H^1(\Sigma)
    \leq l.
  \end{cases}
  \tag{ADP}
\end{equation}

\begin{lemma}\label{lemma:basicinequality}
  Fix $p \geq q > 0$. Then, for all $a, b \geq 0$, we have
  \[
    \frac{p}{q} (a^q - b^q)b^{p-q}\leq a^p - b^p \leq \frac{p}{q}(a^q
    - b^q) a^{p-q}.
  \]
  Here we take the convention that $0^0=1$ when $a=0$ or $b=0$.
\end{lemma}

\begin{proposition}[The ``gradient'' of $\mathscr{J}_p$]
  \label{strongbarycentreapproximation}
  Let $\Sigma \in \mathcal{S}_l$ and let $\xi : \Sigma \to
  \mathbb{R}^d$ be continuous. For $\epsilon > 0$, define
  \begin{equation}\label{eq:perturbedsigma}
    \Sigma_{\epsilon, \xi} \coloneqq \{\sigma + \epsilon \xi(\sigma) \
    | \ \sigma \in \Sigma\} = (\mathbf{1} + \epsilon \xi)(\Sigma).
  \end{equation}
  Suppose $p > 1$ or that $p = 1$ and $\mu(\Sigma) = 0$.
  Then
  \begin{equation} \lim_{\varepsilon \to 0^{+}}
    \frac{\mathscr{J}_p(\Sigma_{\varepsilon, \xi}) -
      \mathscr{J}_p(\Sigma)}{\varepsilon} \leq \max_{\pi_{\Sigma}\in
      \Pi_{\Sigma}}- \int_\Sigma \xi(\sigma) \cdot
    \mathcal{B}_{\pi_\Sigma}(\sigma) \,
    d\nu_{\pi_{\Sigma}}(\sigma). \label{eq:fbary-grad-J-min}
  \end{equation}
\end{proposition}

\begin{definition}\label{def:trivialbarycentrefield}
  Let $\pi_{\Sigma} \in \Pi_{\Sigma}$, and let $\nu_{\pi_{\Sigma}}=
  (\pi_{\Sigma})_{\#}\mu$. We say $\pi_{\Sigma}$ has \textit{trivial}
  barycentre field if
  \[
    \nu_{\pi_{\Sigma}}(\{\sigma \in \Sigma \mid
    \mathcal{B}_{\pi_{\Sigma}}(\sigma) \ne 0\}) = 0.
  \]
  Otherwise, we say the barycentre field of $\pi_{\Sigma}$ is
  \textit{nontrivial}.
\end{definition}

\begin{proposition}[Approximation of
  $\mathcal{B}_{\pi_{\Sigma}}$]\label{approximatebysmooth}
  Let $p \geq 1$. Suppose that $\pi_{\Sigma} \in \Pi_\Sigma$ has
  nontrivial barycentre field. Then, there exists a Lipschitz map $\xi
  : \mathbb{R}^d \to \mathbb{R}^d$ such that
  \[
    \int_{\Sigma} \xi(\sigma) \cdot
    \mathcal{B}_{\pi_{\Sigma}}(\sigma) d \nu_{\pi_{\Sigma}}(\sigma) >
    \frac{1}{2}\int_{\Sigma}|\mathcal{B}_{\pi_{\Sigma}}(\sigma) |^2 d
    \nu_{\pi_{\Sigma}}(\sigma) > 0.
  \]
\end{proposition}

\begin{proposition}[Negligibility of the ambiguous
  locus]\label{ridgesetnegligable}
  Suppose $p > 1$, or $p=1$ with the extra hypothesis $\mu(\Sigma) =
  0$, and let $\Sigma \in \mathcal{S}_l$ be a minimizer of the
  \eqref{eq:adp}. Then $\mu(\mathcal{A}_{\Sigma}) = 0$.
\end{proposition}

\begin{lemma}[Noncut-neighbourhood lemma]\label{noncutneighbourhood}
  Let $\Sigma$ be a locally connected metric continuum (i.e.\ a
  compact, connected, and locally-connected metric space) containing
  more than one point, and let $\sigma \in \Sigma$ be a noncut point
  of $\Sigma$. Then, there exists a sequence $\{B_n\}_{n \in
    \mathbb{N}}$ of open subsets of $\Sigma$ satisfying the following
  conditions:
  \begin{itemize}
    \item[(i)] $\sigma \in B_n$ for all sufficiently large $n$,
    \item[(ii)] $\Sigma \setminus B_n$ is connected for each $n \in
      \mathbb{N}$,
    \item[(iii)] $\diam(B_n) \to 0$ as $n \to \infty$, and
    \item[(iv)] $B_n$ is connected for every $n$.
  \end{itemize}
\end{lemma}

\begin{theorem}[Bounding the mass of noncut
  points]\label{prop:boundingmassofnoncutpoints} Suppose $l > 0$.
  Let $p \geq 2$ and suppose that $\Sigma \in \mathcal{S}_l$ is an
  optimizer and that $\Sigma$ contains at least two points. Let
  $\pi_\Sigma \in \Pi_\Sigma$ and $\nu = (\pi_{\Sigma})_{\#}\mu$.
  Then, there exists some constant $\lambda > 0$ depending only on
  $\mathcal B_{\pi_\Sigma}$ such that for all noncut points $\sigma^*
  \in \Sigma$ we have
  \[
    \nu\{\sigma^*\}|\mathcal{B}_{\pi_{\Sigma}}(\sigma^*)| \geq
    \frac{\lambda }{4l} \int_{\Sigma}
    |\mathcal{B}_{\pi_{\Sigma}}(\sigma)|^2 d\nu(\sigma).
  \]
\end{theorem}

\begin{proof}
  Observe that if $\mathcal{B}_{\pi_\Sigma}(\sigma)$ is trivial
  (\cref{def:trivialbarycentrefield}) then the claim is trivial as
  well; hence suppose $\mathcal{B}_{\pi_\Sigma}(\sigma)$ is
  nontrivial.

  Let $\{B_n\}_{n \in \mathbb{N}}$ be as in
  \cref{noncutneighbourhood}. For each $n \in \mathbb{N}$, let
  $\epsilon_n = \frac{1}{2}\diam(B_{n})$. Define $\Sigma_n = \Sigma
  \setminus B_n$ and let $P_n(x)$ denote the closest-point projection
  onto $\Sigma \cap \partial B_n$. We now bound $\mathscr
  J_p(\Sigma_n)$ by decomposing the domain of integration into two
  pieces and bounding them separately. The first will be
  $\pi_\Sigma^{-1}(\Sigma \setminus \overline{B_n})$, on which by
  construction for all $x \in \pi_\Sigma^{-1}(\Sigma \setminus
  \overline{B_n})$ we have $d(x, \Sigma_n) = d(x, \Sigma)$. The second
  piece will be $\pi_\Sigma^{-1}(\overline{B_n})$, on which by
  construction we have $d(x, \Sigma_n) = d(x, P_n(x))$. In general $|x
  - P_n(x)| \leq |x - P_n(\pi_\Sigma(x))|$, so
  \begin{align*}
    \mathscr{J}_p(\Sigma_n)
    & = \int_{\pi_{\Sigma}^{-1}(\Sigma \setminus
      \overline{B_n})}\dist(x, \Sigma)^p d\mu(x) +
      \int_{\pi_{\Sigma}^{-1}(\overline{B_n})}|x- P_n(x)|^p d\mu(x)\\
    &\leq \int_{\pi_{\Sigma}^{-1}(\Sigma \setminus
      \overline{B_n})}\dist(x, \Sigma)^p d\mu(x) \\
    & \qquad + \int_{\pi_{\Sigma}^{-1}(\overline{B_n})}|x-
      \pi_{\Sigma}(x) - (P_n(\pi_{\Sigma}(x)) - \pi_{\Sigma}(x))|^p
      d\mu(x).
  \end{align*}
  For the last term, since $p \ge 2$, applying
  \cref{lemma:basicinequality} gives
  \begin{align*}
    & |x - \pi_{\Sigma}(x)
      -(P_{n}(\pi_{\Sigma}(x)) - \pi_{\Sigma}(x))|^p - |x -
      \pi_{\Sigma}(x)|^p \\
    &\leq \frac{p}{2}( |x - \pi_{\Sigma}(x) -(P_{n}(\pi_{\Sigma}(x))
      - \pi_{\Sigma}(x))|^2 - |x - \pi_{\Sigma}(x)|^2
      )|x-\pi_{\Sigma}(x)|^{p-2} \\
    &= \frac{p}{2}\Big(
      |P_{n}(\pi_{\Sigma}(x)) - \pi_{\Sigma}(x)|^2
      - 2(P_{n}(\pi_{\Sigma}(x)) - \pi_{\Sigma}(x))\cdot
      (x - \pi_{\Sigma}(x)) \Big)|x - \pi_{\Sigma}(x)|^{p-2},
  \end{align*}
  and since for all $x \in \pi_\Sigma^{-1}( \overline{B_n})$ we have
  $|P_{n}(\pi_{\Sigma}(x)) - \pi_{\Sigma}(x)| \leq \epsilon_n$,
  defining $M = \diam(\supp(\mu))$ and recalling the definition of the
  barycentre field $\mathcal{B}_{\pi_{\Sigma}}$ yields
  \begin{align}\label{eq:Sigma-n-Sigma} \nonumber
    \mathscr{J}_p(\Sigma_{n}) \leq
    & \mathscr{J}_p(\Sigma) - p\int_{\pi_{\Sigma}^{-1}(
      \overline{B_n})}(P_{n}(\pi_{\Sigma}(x)) - \pi_{\Sigma}(x))\cdot
      (x - \pi_{\Sigma}(x))|x - \pi_{\Sigma}(x)|^{p-2} d\mu(x) \\
    \nonumber
    &+ \frac{p}{2} \epsilon_n^2 M^{p-2} \nu(\overline{B_n}) \\
    \nonumber
    =& \mathscr{J}_p(\Sigma) - \int_{ \overline{B_n}}(P_{n}(\sigma) -
       \sigma)\cdot \mathcal{B}_{\pi_{\Sigma}}(\sigma) d\nu(\sigma) +
       \frac{p}{2} \epsilon^2 M^{p-2}\nu(\overline{B_n}) \\
    \leq
    & \mathscr{J}_p(\Sigma) + \epsilon_n
      \int_{\overline{B_n}}|\mathcal{B}_{\pi_{\Sigma}}(\sigma)|
      d\nu(\sigma) +\frac{p}{2} \epsilon_n^2
      M^{p-2}\nu(\overline{B_n}).
  \end{align}

  On the other hand, by \cref{approximatebysmooth} there exists a
  Lipschitz map $\xi: \Sigma \to \mathbb{R}^d$ such that
  \begin{equation}
    \int_{\Sigma}\xi(\sigma)\cdot
    \mathcal{B}_{\pi_{\Sigma}}(\sigma)d\nu(\sigma) >
    \frac{1}{2}\int_{\Sigma}|\mathcal{B}_{\pi_{\Sigma}}(\sigma)|^2
    d\nu(\sigma) > 0. \label{eq:applying-lem-2-11}
  \end{equation}
  Let $L >0$ be a Lipschitz constant for $\xi$, and let $\lambda =
  \max\{\frac{1}{L}, \frac{1}{\max |\xi|}\}$; note $\lambda$ does not
  depend on $\sigma^*$. Then $\lambda \xi$ is $1$-Lipschitz, and thus
  $\sigma \mapsto \sigma + \frac{\epsilon_n}{(l - \epsilon_n)}\lambda
  \xi(\sigma)$ is $1 + \frac{\epsilon_n}{l - \epsilon_n}$-Lipschitz.
  So,
  \[
    \Sigma_n' \coloneqq \{\sigma + \frac{\epsilon_n}{1 -
      \epsilon_n}\lambda \xi(\sigma) \mid \sigma \in \Sigma_{n}\} \in
    \mathcal{S}_l.
  \]

  Now, we want to estimate $\mathscr{J}_p(\Sigma_n') -
  \mathscr{J}_p(\Sigma_n)$ in terms of $\int_{\Sigma}\xi(\sigma)\cdot
  \mathcal{B}_{\pi_{\Sigma}}(\sigma)d\nu(\sigma)$. To do this, we
  first will estimate the difference
  \[
    \bigg|\int_{\Sigma_n} \xi(\sigma)\cdot
    \mathcal{B}_{\pi_{\Sigma_n}}(\sigma)d\nu_{\pi_{\Sigma_n}}(\sigma)
    - \int_{\Sigma \setminus \{\sigma^*\}}\xi(\sigma)\cdot
    \mathcal{B}_{\pi_{\Sigma}}d\nu(\sigma)\bigg|,
  \]
  where we recall that $\nu_{\pi_{\Sigma_n}} \coloneqq
  (\pi_{\Sigma_n})_{\#}\mu$.

  By \cref{ridgesetnegligable} we have $\mu(\mathcal{A}_{\Sigma}) =
  0$, and so for $\mu$-a.e.\ $x\in\mathbb{R}^d \setminus
  \pi_{\Sigma}^{-1}(\overline{B_n})$ we get $\pi_{\Sigma}(x) =
  \pi_{\Sigma_n}(x)$. For the $x \in \pi_\Sigma^{-1}(\overline{B_n})$,
  note $\Sigma_n \subseteq \Sigma$ gives the uniform bound
  \[
    |\pi_\Sigma(x) - \pi_{\Sigma_n}(x)| \leq \mathrm{diam}(\Sigma)
    \leq l.
  \]
  Next, note that regardless of the choices of $\Sigma$ and $x$ we get
  $|\mathcal B_{\pi_\Sigma}(x)| \leq p M^{p-1}$. So the Lipschitz
  condition on $\xi$ gives
  \begin{align*}
    & \bigg|\int_{\Sigma_n} \xi(\sigma) \cdot
      \mathcal{B}_{\pi_{\Sigma_n}}(\sigma)d\nu_{\pi_{\Sigma_n}}(\sigma)
      - \int_{\Sigma \setminus \{\sigma^*\}} \xi(\sigma)\cdot
      \mathcal{B}_{\pi_{\Sigma}} d\nu(\sigma)\bigg| \\
    &\leq pLM^{p-1}\int_{\mathbb{R}^d \setminus \{\sigma^*\}}
      |\pi_{\Sigma}(x) - \pi_{\Sigma_n}(x)| d\mu(x)  \\
    &\leq pLM^{p-1}l\nu(\overline{B_n}\setminus \{\sigma^*\}),
  \end{align*}
  which is $o(1)$. Thus, in particular,
  \begin{equation}
    \varepsilon_n \bigg|\int_{\Sigma_n} \xi(\sigma) \cdot
    \mathcal{B}_{\pi_{\Sigma_n}}(\sigma)d\nu_{\pi_{\Sigma_n}}(\sigma)
    - \int_{\Sigma \setminus \{\sigma^*\}} \xi(\sigma)\cdot
    \mathcal{B}_{\pi_{\Sigma}} d\nu(\sigma)\bigg|  =
    o(\epsilon_n). \label{eq:bary-diff-bound}
  \end{equation}
  Now, by \cref{strongbarycentreapproximation} we have
  \begin{align}
    \mathscr{J}_p(\Sigma_n') - \mathscr{J}_p(\Sigma_n)
    &\leq -\frac{\epsilon_n}{l-\epsilon_n} \int_{\Sigma_n}
      \lambda\xi(\sigma)\cdot
      \mathcal{B}_{\pi_{\Sigma_n}}(\sigma)d\nu_{\pi_{\Sigma_n}}(\sigma)
      + o(\epsilon_n),
      \shortintertext{whence \eqref{eq:bary-diff-bound} gives}
      \mathscr{J}_p(\Sigma_n') -
      \mathscr{J}_p(\Sigma_n)
    &\leq -\frac{\epsilon_n}{l-\epsilon_n} \int_{\Sigma \setminus
      \{\sigma^*\}} \lambda\xi(\sigma)\cdot \mathcal{B}_{\pi_{\Sigma}}
      (\sigma) d\nu(\sigma) + o(\epsilon_n) \nonumber \\
    & = -\frac{\epsilon_n}{l-\epsilon_n} \int_{\Sigma}
      \lambda\xi(\sigma)\cdot \mathcal{B}_{\pi_{\Sigma}} (\sigma)
      d\nu(\sigma) \nonumber \\
    &\qquad + \epsilon_n \nu(\{\sigma^*\}) \lambda \xi(\sigma^*) \cdot
      \mathcal B_{\pi_\Sigma}(\sigma^*) + o(\epsilon_n) \nonumber \\
    &\leq -\frac{\epsilon_n}{l-\epsilon_n}\int_\Sigma \lambda
      \xi(\sigma) \cdot \mathcal B_{\pi_\Sigma}(\sigma) d\nu(\sigma)
      \nonumber \\
    &\qquad + \epsilon_n\nu(\{\sigma^*\}) |\mathcal
      B_{\pi_\Sigma}(\sigma^*)| +
      o(\epsilon_n). \label{eq:Sigma-prime-Sigma-n}
  \end{align}
  Note that the error term in \eqref{eq:Sigma-n-Sigma} is
  $o(\epsilon_n^2)$. So, adding \eqref{eq:Sigma-n-Sigma},
  \eqref{eq:Sigma-prime-Sigma-n} and then applying
  \eqref{eq:applying-lem-2-11} we have
  \begin{align*}
    \mathscr{J}_p(\Sigma_n') - \mathscr{J}_p(\Sigma)
    &< \epsilon_n \bigg(\int_{\overline{B_n}}
      |\mathcal{B}_{\pi_{\Sigma}}(\sigma)|d\nu(\sigma) -
      \frac{1}{2(l-\epsilon_n)}\lambda \int_{\Sigma \setminus
      \{\sigma^*\}} |\mathcal{B}_{\pi_{\Sigma}}(\sigma)|^2
      d\nu(\sigma) \\
    &\quad \qquad + \nu(\{\sigma^*\}) |\mathcal
      B_{\pi_\Sigma}(\sigma^*)| \bigg) + o(\epsilon_n).
  \end{align*}
  Since $\Sigma$ was assumed to be optimal, $\mathscr{J}_p(\Sigma_n') -
  \mathscr{J}_p(\Sigma) \geq 0$, and so for all sufficiently large $n$
  we have
  \[
    \int_{\overline{B_n}}|\mathcal{B}_{\pi_{\Sigma}}(\sigma)|d\nu(\sigma)
    - \frac{\lambda}{2(l-\epsilon_n)} \int_{\Sigma }
    |\mathcal{B}_{\pi_{\Sigma}}(\sigma)|^2 d\nu(\sigma) +
    \nu(\{\sigma^*\}) |\mathcal B_{\pi_\Sigma}(\sigma^*)| \geq 0.
  \]
  Decomposing the left integral via $\overline{B_n} = \{\sigma^*\}
  \cup (\overline{B_n} \setminus \{\sigma^*\})$ and taking $n \to
  \infty$ thus yields that
  \[
    \nu\{\sigma^*\}|\mathcal{B}_{\pi_{\Sigma}}(\sigma^*)| \geq
    \frac{\lambda }{4l}\int_{\Sigma }
    |\mathcal{B}_{\pi_{\Sigma}}(\sigma)|^2 d\nu(\sigma),
  \]
  as desired.
\end{proof}

\end{document}
